\scanpage{16}
\mspage{15}

On peut le faire par RR si $X$ est loc. une intersection complète (rel $/S$), on trouve

\[
\begin{aligned} &\deleted{\unsure{\chi(L_\delta\otimes F)|\chi(L_\delta)}}\\ \chi_F(\delta)&=\chi(L_\delta\otimes F^{\bullet})=\operatorname{tr}\bigl(\operatorname{ch}(L_\delta\,\deleted{\unsure{\otimes F}}{\unsure{\otimes|\,}}F)\operatorname{Todd}(X)\bigr)\\ &\deleted{\unsure{\varphi}}\qquad=\text{terme constant dans }\underbrace{\operatorname{pr}_{2*}\bigl(\operatorname{pr}_1^{*}(\operatorname{ch}_X(L_\delta\otimes F)\inserted{\operatorname{Todd}(X)}\deleted{\unsure{\operatorname{Todd}}})\,\operatorname{ch}_{X\times B}\unsure{\exp(D)|\exp D}\bigr)}_{(\exp D)\,.[\operatorname{ch}_X(L_\delta\otimes F)\operatorname{Todd}(X)]}\end{aligned}
\]

\[
\varphi_F(\delta)=\text{terme de degré 1 dans }\;(\exp D)\,.[\operatorname{ch}_X(L_\delta\otimes F)\operatorname{Todd}(X)]
\]

\uline{Ex. 1}\qquad $X$ lisse $/S$ de dim relative 1. On aura

\[
\operatorname{ch}_X(\underset{\delta\,\text{de degré }d}{L_\delta}\otimes F)=(1+\underset{c(L_\delta)}{\delta})(\underset{\text{rang } F^{\bullet}}{\unsure{\rho|p}}+\underset{\det^{\prime}(F)}{\gamma})=1+\bigl(\underset{\substack{\text{degré }d\\\text{classe }\delta}}{\delta}+\underset{\substack{\text{degré }c\\\deleted{\text{\unsure{catégorie} }C}}}{\gamma}\bigr)
\]

\[
\operatorname{Todd}(X)=1+\underset{\text{degré }1-g}{\underline{\tfrac12 K}}
\]

\[
\operatorname{ch}_X(L_\delta\otimes F)\otimes\operatorname{Todd}(X)=1+\tau,\qquad\text{avec}\qquad\deg\tau=1-g+d+c
\]

On \unsure{a|--} \deleted{Il faut \unsure{prendre le terme|prendre les termes}}

\[
\boxed{\chi_F(\delta)=1-g+d+c}
\]

D'autre part, $\varphi(\delta)$ n'a pas changé \deleted{\unsure{pour}}\inserted{ici} par rapport au calcul précédent dans \unsure{c)|e)}

\[
\boxed{\varphi_F(\delta)=\varphi(\delta)=\tfrac12\operatorname{pr}_{2*}(D^2)=\text{polarisation canonique de }B}
\]

Donc on trouve une section canonique de

\[
\underline{\mathrm{Alb}}_{X/S}^{\chi_F(\delta)=\chi(\delta)+c}\overset{A}{\times}\mathrm{Pic}_{B/S}^{\varphi(\delta)}\overset{A}{\times}\bigl(\mathrm{Pic}_{X/S}^{\deleted{\unsure{\varphi(}}\delta\deleted{\unsure{)}}}\overset{B}{\times}(B,\widetilde{\varphi(\delta)})\bigr)
\]

donc, en utilisant le résultat fait dans c, on trouve une section canonique de $\underline{\mathrm{Alb}}^{c}_{X/S}\simeq\underline{\mathrm{Pic}}^{c}_{X/S}$.

Il n'y a pas de mérite à cela, car il suffit de prendre la section définie par \inserted{le Module inversible} $\det{}^*F$, sur $X/S$.
